\section{Scope, Method, and Qualifications}\label{sec:scope-method}

\subsection*{Audience, Series Place, and Governing Question}

This companion essay accompanies \emph{The Creature Before God} and addresses intelligent adult readers---believers, skeptics, and students of science or philosophy---who encounter the widespread assumption that physical measurement and empirical verification exhaust what is real or rationally knowable. Its governing question is strictly epistemological and methodological: whether the success of physical instrumentation warrants excluding non-metrical realities and rational metaphysics, and what cognitive standing natural reasoning from sensible effects toward God retains once that exclusion fails.

Because its role is ground-clearing rather than full metaphysical execution, the essay analyzes the logical structure of \emph{demonstratio quia} and the starting points of St.\ Thomas Aquinas's Five Ways without attempting to execute or defend each individual proof at treatise length; full demonstrations of natural theology belong to the main work. Likewise, refuting instrumental exclusivism clears away an illegitimate prior disqualification without claiming that every unmeasurable assertion is thereby proved true.

\subsection*{Evidentiary and Historical-Analysis Boundaries}

Two distinct evidentiary strata govern the citations in this work and are kept rigorously separate:

\begin{itemize}
\item \emph{Locally verified public-domain primary texts.} Every verbatim quotation in the body and endnotes comes directly from a tracked, public-domain witness inspected at its exact locus: Sacred Scripture (Challoner Douay-Rheims Bible, paired with the Clementine Vulgate), St.\ Augustine (\emph{Confessions} X.6 via CCEL Nicene and Post-Nicene Fathers, Series 1, vol.\ 1; \emph{On the Trinity} Book III, ch.\ 3 in Arthur West Haddan's translation, revised by W.\ G.\ T.\ Shedd in NPNF, Series 1, vol.\ 3, 1887, New Advent web transcription, uncollated against print), St.\ Thomas Aquinas (\emph{Summa Theologiae}, Prima Pars, in the English Dominican translation, Project Gutenberg eBook 17611), St.\ John Henry Newman (\emph{Parochial and Plain Sermons}, vol.\ 1, London: Rivingtons, 1869 OCR), G.\ K.\ Chesterton (\emph{Orthodoxy}, 1908, Project Gutenberg eBook 16769), the First Vatican Council (\emph{Dei Filius}, 1870, Latin text and Cardinal Manning's English translation in \emph{Dogmatic Canons and Decrees}, 1912), Pope Leo XIII (\emph{Aeterni Patris}, 1879, Latin text and 1903 Benziger English translation), and Pope St.\ Pius X (\emph{Pascendi Dominici Gregis}, 1907, Latin text and 1908 English translation appended to \emph{The Programme of Modernism}).
\item \emph{Bibliographic and conceptual engagement without verbatim quotation.} Recent magisterial texts under Vatican copyright---Second Vatican Council, \emph{Gaudium et spes} (nos.\ 36, 57); \emph{Catechism of the Catholic Church} (nos.\ 31--39); and Pope St.\ John Paul II, \emph{Fides et ratio} (nos.\ 55, 67, 83, 88)---are bound at the cataloged edition level and presented strictly in original doctrinal paraphrase with zero verbatim reproduction. Modern philosophical, metrological, and popular-scientific figures whose works are not locally tracked in the source library---David Hume, Auguste Comte, Pierre Duhem, the Vienna Circle (Hahn, Neurath, Carnap), A.\ J.\ Ayer, Karl Popper, Alonzo Church, Carl Hempel, W.\ V.\ O.\ Quine, Bas van Fraassen, the JCGM metrological guides (VIM/GUM), Stephen Hawking and Leonard Mlodinow, and Lawrence Krauss---are analyzed accurately in original prose with standard bibliographic citations and zero fabricated or verbatim quotations.
\end{itemize}

\subsection*{Theological Authority and Reliance Qualifications}

Doctrinal authorities remain visibly distinguished by weight: inspired Sacred Scripture (Wis.\ 11:21; 13:1--9; Rom.\ 1:18--25), solemn dogmatic definition of an ecumenical council (First Vatican Council, \emph{Dei Filius}, especially Canons I.2 and II.1), conciliar teaching of an ecumenical council (\emph{Gaudium et spes}), papal encyclicals (\emph{Aeterni Patris}, \emph{Pascendi Dominici Gregis}, \emph{Fides et ratio}), the Church's universal catechetical synthesis (\emph{Catechism of the Catholic Church}), patristic and scholastic doctors (St.\ Augustine, St.\ Thomas Aquinas), Catholic literary and pastoral witnesses (St.\ John Henry Newman, whose \emph{Parochial and Plain Sermons}, vol.\ 1, were preached during his Anglican ministry at St.\ Mary's, Oxford, 1825--1843, and republished unaltered by W.\ J.\ Copeland at Rivingtons in 1868/1869 with Newman's permission, with Copeland's preface explicitly cautioning that republication does not constitute Catholic-era re-assertion of every passage; G.\ K.\ Chesterton), and original philosophical analysis. Offered strictly as a study aid in Catholic philosophy and theology, the material carries no magisterial authority and serves as neither a critical edition nor an official ecclesiastical document.

\subsection*{Rights and Review Boundaries}

Public-domain primary texts and historical translations retain their independent public-domain status; copyrighted magisterial editions and modern scholarly or scientific works are cited bibliographically without reproducing protected expression. Original exposition and analytical prose are licensed under CC BY 4.0 as stated in the colophon. All verbatim quotations were verified directly against tracked local witnesses, and structural, binding, web-edition, and print checks were run locally; no external specialist, clerical, or peer review is claimed.
